\documentclass[exercice,a4]{thedocument}%
\startdatakeys
    \source{laclassecooperativedestef.com}
    \BO{2026}
    \cycle{3}
    \competencesDuSocle{RE4,CO2}
    \objectifsApprentissage{C3C1ab,C3C1aj,C3C1ad,C3C2ac,C3C1ae}
    \motsCles{cercle, droite, perpendicularite, triangle}
\enddatakeys
\begin{document}%
\noindent\textbf{Programme de construction 10}\par\smallskip\noindent
\begin{itemize}
    \item Tracer une droite $(\mathspoint{d})$ et placer le point
    $\mathspoint{A}$ dessus.
    \item Tracer sur la droite $(\mathspoint{d})$ un arc de cercle
    $\mathcal{C}$ de centre $\mathspoint{A}$ et de rayon $3\text{ cm}$. Il
    coupe la droite $(\mathspoint{d})$ en $\mathspoint{B}$ et
    $\mathspoint{C}$.
    \item Tracer le segment $[\mathspoint{AD}]$ à l'opposé de
    $\mathspoint{C}$, sachant que $\mathspoint{AD}=3{,}5\text{ cm}$ et que
    $[\mathspoint{AD}]$ est perpendiculaire à $(\mathspoint{d})$.
    \item Tracer $[\mathspoint{BD}]$ et $[\mathspoint{DC}]$.
\end{itemize}
\end{document}%
